\documentclass[10pt,twocolumn]{article}

\usepackage[margin=0.60in,columnsep=0.22in]{geometry}
\usepackage{amsmath,amssymb,booktabs,enumitem,microtype}
\usepackage[font=small,labelfont=bf]{caption}
\usepackage{xcolor}
\usepackage[colorlinks=true,allcolors=blue!55!black]{hyperref}

\setlength{\parindent}{0pt}
\setlength{\parskip}{3pt}
\setlength{\emergencystretch}{2em}
\setlist{nosep,leftmargin=*}
\newcommand{\D}{\Delta_{60}}
\newcommand{\BR}{\operatorname{BR}}

\title{\vspace{-1.2em}\textbf{Perfect Mode Hal: Aggressive, Change-Aware\\
Opponent Exploitation in Exact DTH}}
\author{James Cenawood}
\date{August 20, 2026}

\begin{document}
\maketitle
\vspace{-1.3em}

\begin{abstract}
Perfect Mode (PM) Hal is a causal meta-policy for repeated, pure Drop the
Handkerchief (DTH). One certified tablebase remains the sole authority for each
continuation-adjusted $60\times60$ matrix. Around it, PM combines Adaptive
Hal's Bayesian counts, Perfect Hal's pattern experts, online change-point and
outcome models, and a compatible Aggro Hal recurrent network. Fixed Share
learns role-specific source and controller weights. Four modes move from an
exact shield to hard forecast responses, while every selected policy receives
a fresh local worst-case-loss measurement. An earlier 24-identity development
run scored 44 wins in 48 games, but its protocol lacked durable registration
and its evaluator trusted self-reported risk. It has no promotion authority.
This note instead registers a 56-identity confirmation protocol with independent
matrix recomputation, a no-Aggro ablation, and family and seat gates. At the
protocol commit that confirmation had not been executed. The one-shot run
scored PM at $.8705$ (97--14--1), behind Perfect at $.9554$; the promotion gate
failed despite zero independent local-risk audit violations. PM is an
aggressive candidate, not human mind reading, dominance, or a whole-game safety
theorem.
\end{abstract}

\section{Scope and design philosophy}

At public state $x$, the exact DTH agent supplies matrix $M(x)$, equilibrium
policies $(x^\star,y^\star)$, and a saddle-gap certificate at most $10^{-6}$.
All actions are literal seconds $1,\ldots,60$. PM never estimates continuation
values and never sees an unrevealed simultaneous action.

The intended character is ``overbearing after evidence, composed after
surprise.'' Perfect contributes hard exploitation; Adaptive contributes
uncertainty and a constrained response frontier; Aggro contributes recurrent
features and a learned residual policy. The story evidence ledger motivates
competing personae, opponent mental models, and abrupt memory discontinuities.
PM translates that fiction into testable engineering---model competition,
role separation, recency, and change recovery---rather than treating a person
as one fixed psychological type. The final narrative ideal, to improve and
adapt without looking back, becomes persistent memory with rapid de-weighting
and recovery after changed play.

\section{Causal opponent model}

For each opponent role separately, PM maintains six mandatory forecast
sources: uniform recovery, the opponent's exact equilibrium policy, a decayed
Dirichlet model, Perfect's online expert mixture, categorical Bayesian online
change-point detection (BOCPD), and an outcome-conditioned model. A seventh
source is Aggro's GRU forecast when a strict compatible checkpoint is present.
After action $a_t$ is revealed, source $k$ receives prequential log score and
Fixed Share updates
\begin{align*}
 \widetilde w_{k,t+1}&\propto w_{k,t}
       \exp\!\left(\eta\log q_{k,t}(a_t)\right),\\
 w_{k,t+1}&=(1-\alpha)\widetilde w_{k,t+1}
 +\frac{\alpha}{K-1}\left(1-\widetilde w_{k,t+1}\right),
\end{align*}
with $\alpha=.04$. Thus every source retains recovery mass when an opponent
switches. BOCPD carries a truncated Dirichlet--multinomial run-length posterior
with hazard $.08$. The fused forecast is
$\widehat q_t=\sum_k w_{k,t}q_{k,t}$.

Aggression confidence is the declared heuristic
\[
 c_t=c_{\max}\frac{n_{\rm eff}}{4+n_{\rm eff}}
 (1-P_t(\text{change}))(1-\mathrm{JS}_t)s_t,
\]
where $\mathrm{JS}_t$ is normalized source disagreement and
$s_t=\exp[-\max(0,\overline{\mathrm{NLL}}_t-\log60)]$. Only the evidence
fraction has a data-biased-response interpretation; the product is a project
controller, not a published theorem.

\section{Aggression with measured risk}

For Dropper, the certified $\epsilon$ candidate solves
\[
 \max_{x\in\D}x^{\mathsf T}M\widehat q_t
 \quad\text{s.t.}\quad
 (M^{\mathsf T}x)_j\ge L-\epsilon\quad\forall j,
\]
where $L=\min_j(M^{\mathsf T}x^\star)_j$ is freshly recomputed. Checker uses
the symmetric upper constraint. The bank also contains exact equilibrium,
hard responses to fused and Perfect forecasts, and Aggro's direct policy. For
a direct Dropper policy $x$, actual local loss is
$\max(0,L-\min_j(M^{\mathsf T}x)_j)$; every direct and frontier candidate is
independently checked against the current matrix.

One common bank of heuristic ensemble perturbations must give positive expected
improvement and at least $.55$ positive-draw support. These draws are not a
calibrated posterior, so the support fraction is not an improvement probability.
Cold start or a change shock enters \emph{shield}, where only the exact face is
admissible. The registered v3 caps are \emph{probe} $.05$, \emph{press} $.50$,
and \emph{dominate} $2.0$, inside a $12.0$ per-game charge budget. Controller
weights receive full-information counterfactual payoff feedback after reveal.
This is deliberately aggressive: the budget is an explicit risk appetite, not
a guarantee that losses add across a DTH game. Ganzfried--Sandholm's cumulative
ledger concerns repetitions of one stage game; DTH's matrices change with the
state and already contain overlapping continuation values.

\section{Development audit and registered confirmation}

The v1 and v2 files were run before this branch registered them in Git; mutable
status text cannot establish preregistration. They remain development evidence.
V2 used 24 identities across six families, common seeded seat pairs, and 10,000
identity-clustered bootstrap replicates. Stopped games scored $.5$.

\begin{table}[h]
\centering
\caption{Development v2 and registered v3 outcomes.}
\small
\begin{tabular}{lrrr}
\toprule
Controller & V2 score & V3 score & V3 W--L--S\\
\midrule
PM Hal     & .9167 & .8705 & 97--14--1\\
Perfect    & .8958 & .9554 & 107--5--0\\
PM no Aggro& ---   & .8259 & 92--19--1\\
Aggro      & .6458 & .7232 & 81--31--0\\
Adaptive   & .6563 & .7188 & 80--31--1\\
Exact      & .5417 & .4955 & 54--55--3\\
\bottomrule
\end{tabular}
\end{table}

V2's nominal every-component gate failed. Its controller-reported risk was not
independently reconstructed, so it cannot promote PM. Fused expected NLL was
$3.481$ $[2.904,3.978]$, but change detection found only 10 of 86 truth shifts.
These mixed results motivated the registered repair.

Clean-Git v3 hash-binds the controller, recurrent checkpoint, and tablebase,
records its commit, and uses 56 fresh identities over all 14 families (112 games
per controller). It compares PM with Exact, Perfect, Adaptive, Aggro, and PM
without Aggro. Promotion requires zero independent risk violations and
nonnegative aggregate, both-seat, and every-family comparisons. Every selected
policy's loss, charge contract, mode cap, cumulative charge, and game budget is
recomputed from its continuation-adjusted matrix.

The registered protocol ran once from commit \texttt{b76147f}.
PM-minus-baseline paired
score differences were $-.0848$ $[-.1518,-.0223]$ against Perfect,
$+.0446$ $[.0089,.0893]$ against PM without Aggro, and $+.3750$
$[.2723,.4777]$ against Exact. The gate failed the every-component, both-seat,
and every-family conditions; Perfect was the strongest component. The
independent audit found zero violations over 788 decisions and 112 games, with
maximum measured local loss $1.5373$. Fused expected NLL was $3.5524$ versus
uniform $4.0943$, but the change detector found only 35 of 168 truth shifts and
produced 23 false-alarm onsets. The result supports aggregate value from enabling
the Aggro component in this protocol; it does not isolate recurrent memory as
the cause. The hash-bound compact record is
\texttt{pm\_hal\_confirmation\_v3.json}.

\section{Human-play interpretation and limits}

The psychology is operational: reinforcement, belief learning, response to
one's own revealed action, recency, periodicity, and change are competing
causal hypotheses. A human can counter-adapt, deceive, or change strategy; PM
must preserve recovery mass and retreat on surprise. It should never label a
participant's personality or claim privileged access to intent.

A human study must keep opponent identity across games, log the sequence
forecast before each reveal, score realized NLL/Brier without simulator-truth
distributions, and cluster outcomes by participant. Change delay and false
alarms, role slices, calibration, and game score are co-primary diagnostics.
Bounded recursive reasoning is one possible population archetype, not a label
the controller may assign permanently to a participant.
The present measured evidence comprises development evidence from 24 identities
and one registered synthetic confirmation over 56 fresh identities and known
families. The recurrent checkpoint saw 2,653 training decisions. Its enabled
component improved aggregate score over the no-Aggro ablation, while a causal
recurrent-memory advantage remains open. PM is feature-complete for human play
as an aggressive candidate. Synthetic promotion failed; any future protocol
requires a new version and fresh identities, followed by participant-level human
evidence.

\begingroup
\scriptsize
\setlength{\parskip}{0pt}
\begin{thebibliography}{9}
\bibitem{rnr} M. Johanson, M. Zinkevich, and M. Bowling.
\href{https://papers.nips.cc/paper_files/paper/2007/file/6e7b33fdea3adc80ebd648fffb665bb8-Paper.pdf}{Computing Robust Counter-Strategies}. NeurIPS, 2007.
\bibitem{dbr} M. Johanson and M. Bowling.
\href{https://proceedings.mlr.press/v5/johanson09a.html}{Data Biased Robust Counter Strategies}. AISTATS, 2009.
\bibitem{share} M. Herbster and M. Warmuth.
\href{https://doi.org/10.1023/A:1007424614876}{Tracking the Best Expert}. Machine Learning, 1998.
\bibitem{bocpd} R. Adams and D. MacKay.
\href{https://arxiv.org/abs/0710.3742}{Bayesian Online Changepoint Detection}, 2007.
\bibitem{ewa} C. Camerer and T. Ho.
\href{https://doi.org/10.1111/1468-0262.00054}{Experience-Weighted Attraction Learning}. Econometrica, 1999.
\bibitem{mentalize} A. Hampton, P. Bossaerts, and J. O'Doherty.
\href{https://doi.org/10.1073/pnas.0711099105}{Mentalizing-related computations in strategic interaction}. PNAS, 2008.
\bibitem{delta} M. Nassar et al.
\href{https://doi.org/10.1523/JNEUROSCI.0822-10.2010}{A Bayesian delta-rule model of changing environments}. J. Neuroscience, 2010.
\bibitem{ch} C. Camerer, T.-H. Ho, and J.-K. Chong.
\href{https://doi.org/10.1162/0033553041502225}{A Cognitive Hierarchy Model of Games}. QJE, 2004.
\bibitem{scores} T. Gneiting and A. Raftery.
\href{https://doi.org/10.1198/016214506000001437}{Strictly Proper Scoring Rules}. JASA, 2007.
\bibitem{safe} S. Ganzfried and T. Sandholm.
\href{https://doi.org/10.1145/2716322}{Safe Opponent Exploitation}. ACM TEAC, 2015.
\end{thebibliography}
\endgroup

\end{document}
